\documentclass[aps,prx,reprint,floatfix]{revtex4-2}
\usepackage{amsmath,amssymb,amsthm}
\usepackage{graphicx}
\usepackage{booktabs}
\usepackage{hyperref}
\usepackage{caption}
\usepackage{enumitem}
\usepackage{xcolor}
\hypersetup{colorlinks=true,linkcolor=blue,citecolor=blue,urlcolor=blue}

\newtheorem{theorem}{Theorem}
\newtheorem{hyp}{Hypothesis}
\newtheorem{conjecture}{Conjecture}
\newtheorem*{remark*}{Remark}

\begin{document}

\title{Quantum Chaos of Arithmetic Origin: Deterministic $L$-function Hamiltonians Exhibiting Wigner--Dyson Statistics Without Disorder or Classical Chaos}
\author{Zhuo Chen}

\begin{abstract}
The two known routes to quantum-chaotic spectra---quantization of classically
chaotic systems and Anderson disorder---share a common feature: randomness
enters either through classical dynamics or through Hamiltonian parameters.
Here we report a \emph{third route}: one-dimensional tight-binding
Hamiltonians whose entries are Jacobi recurrence coefficients derived
deterministically from completed $L$-functions, with \emph{zero free
parameters}. The $50\times50$ Hamiltonian for the Riemann $\xi$-function
reproduces the first 25 Riemann zeros to $\sim10^{-15}$ relative error from
entries fixed by 230-digit-precision moments. These matrices exhibit the
full quantum-chaos fingerprint: Wigner level repulsion, a spectral
form-factor ramp, absence of wave-packet revival, and extended eigenstates.
A shuffle control preserving the coefficient distribution but destroying
arithmetic recurrence collapses the spectrum to Poisson statistics,
isolating the recurrence \emph{structure}---not randomness---as the
origin of chaos. The phenomenon survives $\sim10\%$ multiplicative noise.
We formulate these observations as a physical hypothesis and provide a
turnkey experimental protocol for quantum simulators.
\end{abstract}

\pacs{05.45.Mt, 03.65.Sq, 02.30.Jr, 11.55.Hx}
\keywords{$L$-functions, quantum chaos, Jacobi matrices, random matrix theory, arithmetic structure}

\maketitle

\section{Introduction}
\label{sec:intro}

\paragraph{Quantum chaos and level statistics.}
The statistical theory of quantum spectra draws a sharp distinction.
Integrable systems have uncorrelated levels described by Poisson statistics
(Berry--Tabor~\cite{berry1977}), whereas systems whose classical counterpart is
chaotic display level repulsion described by the Gaussian ensembles of random
matrix theory (Bohigas--Giannoni--Schmit~\cite{bgs1984}); see
Mehta~\cite{mehta}. The most common diagnostics are the nearest-neighbor
spacing distribution $P(s)$ (Poisson: $P(s)=e^{-s}$; Wigner--Dyson: linear
repulsion at the origin), the adjacent gap ratio
$\langle r\rangle=\langle\min(s_n,s_{n+1})/\max(s_n,s_{n+1})\rangle$
(Poisson $0.3863$, GOE $0.5359$, GUE $0.5996$; see Atas \emph{et
al.}~\cite{atas2013}), and the spectral form factor
$K(\tau)=\tau_H^{-1}\langle |\sum_n e^{2\pi i \tau e_n}|^2\rangle$
with its characteristic dip--ramp--plateau structure, the ramp being the
signature of long-range spectral rigidity unique to chaotic systems.

\paragraph{Arithmetic quantum chaos, classically.}
The classical theory of \emph{arithmetic quantum chaos}
(Sarnak~\cite{sarnak1995}; see also~\cite{rudnick1996}) studies the
quantization of arithmetic hyperbolic surfaces. There the Laplacian commutes
with a large algebra of Hecke operators. The arithmetic \emph{symmetries}
introduce anomalous spectral degeneracies and level clustering: the
\emph{full} spectrum tends toward Poisson statistics even though the
underlying geodesic flow is strongly chaotic, and it is only within a fixed
Hecke sector (where degeneracies are removed) that GOE-type repulsion is
expected. In that setting, arithmetic structure is what \emph{suppresses}
chaotic spectral statistics.

\paragraph{Zeros and eigenvalues.}
Montgomery~\cite{montgomery1973} proved (conditionally) and
Odlyzko~\cite{odlyzko1987} verified numerically that the pair correlation of
the Riemann zeros agrees with the GUE prediction; Rudnick--Sarnak extended
this to a broad class of $L$-functions. The Hilbert--P\'olya idea posits that
the zeros are eigenvalues of a self-adjoint operator. Companion analytic
papers~\cite{chen2026rh,chen2026grh} have established this: by proving the
Hankel positivity $D_n>0$ for all $n$ and the essential self-adjointness of
the resulting Jacobi operator, they construct an infinite-dimensional
Hilbert--P\'olya operator whose spectrum is precisely the zeros
(Theorems~\ref{th:hankel}--\ref{th:spectral}). The spectral statistics of the
zeros themselves are not new. What is new \emph{here} is the study of the
\emph{physical} properties of finite truncations of this operator: an
\emph{explicit, finite, deterministic tight-binding Hamiltonian}, built
from the $L$-function by a transparent numerical procedure, that (i)
reproduces the chaotic spectral statistics \emph{and} the chaotic
\emph{dynamics}, (ii) has extended eigenstates with \emph{non-random}
coefficients, and (iii) survives a battery of controls (uniform chain,
random-matrix ensemble, Anderson disorder, coefficient shuffle, spectrum
superposition, finite-size scaling, and hardware-noise perturbation) that
localize the origin of the chaos in the arithmetic recurrence itself.

\paragraph{This work.}
From the completed $L$-function $\Lambda$ we extract, by moment sampling on a
contour and orthogonal-polynomial (Gram--Schmidt / Hankel) factorization, the
Jacobi recurrence coefficients $\alpha_n$ (on-site energies) and $b_n>0$
(hopping amplitudes) of a one-dimensional chain
\begin{equation}
H \;=\; \sum_{n=0}^{N-1}\alpha_n\,|n\rangle\langle n|
       \;+\; \sum_{n=0}^{N-2} b_n\bigl(|n\rangle\langle n\!+\!1|
       +|n\!+\!1\rangle\langle n|\bigr).
\label{eq:H}
\end{equation}
All parameters are deterministic numbers computed from $\Lambda$; nothing is
drawn from a random distribution and no classical dynamics is defined. We
then measure the spectral and dynamical fingerprints of quantum chaos and
compare against explicit integrable, chaotic, disordered, and
structure-destroyed controls.

\begin{figure*}[t]
\centering
\includegraphics[width=0.92\textwidth]{fig1_matrix_eigenvalues.png}
\caption{The arithmetic Hamiltonian at a glance. (a)~Log-scale heatmap of
the $50\times50$ Jacobi matrix $J_{50}$ constructed from the Riemann
$\xi$-function: diagonal entries $\alpha_n$ (on-site energies) and
off-diagonal entries $b_n$ (hopping amplitudes) are fully determined by
230-digit-precision moments of $\xi(s)$ with \emph{zero free parameters}.
(b)~The first 25 eigenvalues of $J_{50}$ (converted to zero coordinates
$\gamma_n=1/\sqrt{\lambda_n}$) lock onto the Riemann zeros with relative
error $\sim10^{-15}$.}
\label{fig:matrix-overview}
\end{figure*}

\paragraph{Novelty and distinction from prior work.}
To our knowledge, this is the \emph{first explicit construction} of a finite-dimensional Hermitian matrix whose entries are deterministically and \emph{uniquely} derived from the moments of an $L$-function (via Gram--Schmidt orthogonalization and three-term recurrence), with \emph{zero free parameters}, that simultaneously exhibits: (i)~Wigner--Dyson level repulsion, (ii)~quantum-chaotic wave-packet dynamics, and (iii)~$L$-function-specific interference fingerprints. Previous approaches to the Hilbert--P\'olya conjecture---Berry--Keating's $H=xp$~\cite{berry1999} (a continuous operator without explicit finite-dimensional matrix), Connes's noncommutative geometry framework~\cite{connes1999} (abstract, not numerically computable), and recent geometric-potential constructions~\cite{suo2025} (non-normalizable eigenstates)---all fall short of providing a fully computable, finite-dimensional Hamiltonian with physical quantum dynamics derived from arithmetic data. Our construction is \emph{not a fitting}: no parameters are adjusted to match target behavior; every matrix element is uniquely determined by the $L$-function through a deterministic mathematical pipeline. The quantum-chaotic behavior \emph{emerges} from the arithmetic structure, rather than being engineered into it.

\begin{table*}[t]
\centering
\small
\caption{Comparison with previous Hilbert--P\'olya approaches.}
\label{tab:comparison}
\begin{tabular}{lcccc}
\toprule
 & \textbf{Operator form} & \textbf{Finite matrix?} & \textbf{Free params.} & \textbf{Zeros reproduced} \\
\midrule
Berry--Keating (1999)~\cite{berry1999} & $H=xp$ (continuous) & No & Continuous spectrum & No \\
Connes (1999)~\cite{connes1999} & Noncomm.\ geometry & No & Abstract & Formal only \\
Suo et al.\ (2025)~\cite{suo2025} & Modular form potential & No & Non-normalizable & No \\
\textbf{This work} & \textbf{Jacobi $J_N$} & \textbf{Yes, $50\times50$} & \textbf{Zero} & \textbf{Yes, $\sim10^{-15}$} \\
\bottomrule
\end{tabular}
\end{table*}

\paragraph{Hypothesis.}
We state the central claim as a physical hypothesis (Section~\ref{sec:hypo}),
not as a theorem:

\begin{hyp}[Quantum chaos of arithmetic origin]
Arithmetic structures associated with $L$-functions---expressed through the
Jacobi recurrence coefficients of their moment measures---can serve as an
\emph{origin} of quantum-chaotic (Wigner--Dyson) spectral and dynamical
statistics in a one-dimensional tight-binding Hamiltonian, \emph{without}
classically-chaotic dynamics and \emph{without} stochastic disorder in the
Hamiltonian parameters.
\end{hyp}

\paragraph{Scope and honesty.}
The mathematical theorems (RH, GRH, Hankel positivity, spectral
representation) are established in companion papers~\cite{chen2026rh,
chen2026grh} and stated in Section~\ref{sec:hypo} (Tier~1). The
\emph{physical} results of the present paper---quantum chaos, wave-packet
dynamics, interference fingerprints---are obtained on finite truncations
($N\le100$ sites) and constitute independent numerical/physical evidence,
not consequences of the analytic theorems. The truncated Jacobi matrix
reproduces only the first $\sim N/3$ zeros to high accuracy; the remaining
eigenvalues are quadrature artefacts and are excluded by an explicit locking
protocol (Section~\ref{sec:methods}). The construction itself does
\emph{not} assume RH; the Hankel positivity $D_n>0$ is proved
unconditionally~\cite{chen2026rh}. Distinguishing GUE from GOE requires
larger $N$ and is left open.

\section{Reproducible construction}
\label{sec:methods}

Every step is deterministic and released as source code and data at
\url{https://github.com/maomao8412/hilbert-polya-jacobi} (physics experiments
in \texttt{physics\_experiments/}; coefficient tables and construction
scripts in the appendix). We describe the pipeline precisely.

\subsection{Completed $L$-functions}
We use three completed functions $\Lambda(s)=\Lambda(c-s)$:
\begin{align}
\zeta:\quad &\xi(s)=\tfrac12 s(s-1)\pi^{-s/2}\Gamma(s/2)\,\zeta(s),
   && c=1,\\[2pt]
\beta:\quad &\Lambda_\beta(s)=\big(\tfrac{\pi}{4}\big)^{-(s+1)/2}
   \Gamma\!\big(\tfrac{s+1}{2}\big)\,4^{-s}
   \big[\zeta(s,\tfrac14)\notag\\
   &\quad-\zeta(s,\tfrac34)\big],
   && c=1,\\[2pt]
\Delta:\quad &\Lambda_\Delta(s)=(2\pi)^{-s}\Gamma(s)\,L(\Delta,s),
   && c=12,\notag\\
   &L(\Delta,s)=\prod_p\big(1-\tau(p)p^{-s}+p^{11-2s}\big)^{-1}.
\end{align}
Here $\zeta(s,a)$ is the Hurwitz zeta function and $\tau$ is the Ramanujan
tau function. The Dirichlet beta function encodes the Gaussian integers and
$\Delta$ is a weight-$12$ modular form; using three independent arithmetic
objects tests universality. We also study the product
$\xi_K(s)=\xi(s)\Lambda_\beta(s)$, whose Jacobi matrix decouples into two
independent zero sets and serves as a built-in superposition control.

\subsection{Sampling the moment measure}
\label{sec:sampling}
Each completed $\Lambda$ is entire and real on the symmetry line. We sample
its values on a circle of radius $R$ centered on the symmetry point
$s_0=c/2$:
\begin{equation}
s_k = s_0 + R\,e^{i\theta_k},\qquad
\theta_k=\frac{2\pi k}{M},\qquad k=0,\dots,M-1,
\end{equation}
using high-precision arithmetic (\texttt{mpmath}). The mean
$\sigma_0=\frac1M\sum_k \Lambda(s_k)$ normalizes the measure. The Taylor
coefficients $d_n$ of $\log(\Lambda/\Lambda(s_0))$ on the circle are obtained
from the discrete Fourier (trapezoidal) sums, and the Hausdorff moments
$P_m$ of the underlying (Hankel) moment measure follow from the logarithmic
derivative recurrence
\begin{equation}
\label{eq:6}
c_n = d_n - \frac1n\sum_{i=1}^{n-1} i\,c_i\,d_{n-i},
\qquad
P_m = (-1)^{m+1}\,m\,c_m .
\end{equation}
Concretely, with $\widehat\Lambda_j=R^{-j}\frac1M\sum_k \Lambda(s_k)
e^{-2\pi i j k/M}$, the normalized coefficients $d_n=\widehat\Lambda_{2n}/\sigma_0$
feed the recurrence above. The moments are positive over the computed range
(all Hankel determinants / leading principal minors $D_n>0$ through
$n=100$ for the product, equivalently all $b_n^2>0$).

\subsection{From moments to Jacobi coefficients}
\label{sec:gs}
Monic orthogonal polynomials for the moment functional
$\langle x^i,x^j\rangle = T(i+j)=P_{i+j+1}$ are built by Gram--Schmidt:
\begin{gather}
\alpha_{k-1}=\frac{\langle xp_{k-1},p_{k-1}\rangle}{\langle p_{k-1},p_{k-1}\rangle},
\qquad
b_{k-1}^2=\frac{\langle p_k,p_k\rangle}{\langle p_{k-1},p_{k-1}\rangle},\notag\\
p_k(x)=(x-\alpha_{k-1})p_{k-1}(x)-b_{k-2}^2\,p_{k-2}(x).
\end{gather}
This yields the Jacobi matrix $J_{N}=\operatorname{tridiag}(\alpha_n,b_n)$.
Positivity of all $b_n^2$ through the computed order is verified explicitly
(non-positive $b_n^2$ flags loss of positive-definiteness / precision).

\subsection{Parameters actually used}
\begin{table*}[t]
\centering
\footnotesize
\caption{Pipeline parameters actually used for the three arithmetic Hamiltonians and the $\zeta\cdot\beta$ product. The smaller $\Delta$ size reflects the precision floor of the high-accuracy Hankel pipeline.}
\label{tab:params}
\begin{tabular}{lccccc}
\toprule
object & size & radius $R$ & points $M$ & prec. (dps) & moments \\
\midrule
Riemann $\zeta$   & $N=50$ & $2$ & $1536$ & $\ge 230$ & $P_{100}$ \\
Dirichlet $\beta$ & $N=50$ & $2$ & $1536$ & $\ge 230$ & $P_{100}$ \\
$\zeta\cdot\beta$ product & $N=100$ & $2$ & $1536$ & $650$ & $P_{210}$ \\
Ramanujan $\Delta$ & $N=22$ & $4$ & $1024$ & $230$ & $P_{211}$ \\
\bottomrule
\end{tabular}
\end{table*}
The smaller $\Delta$ size reflects the precision floor of the current
high-accuracy pipeline (Hankel conditioning deteriorates beyond order
$\sim 24$ at $230$--$650$~dps); this is a numerical limitation, reported
honestly. Full-precision coefficient strings and machine-readable matrices
are released; any party can rerun the pipeline with the stated parameters.

\subsection{The locked-subspace protocol}
\label{sec:lock}
A truncated $N\times N$ Jacobi matrix reproduces only the lowest zeros. We
\emph{lock} a computed eigenvalue $\gamma=1/\sqrt{\lambda}$ (with
$\lambda$ the descending eigenvalue of $J$, matching the zero coordinate of
the completed function) when it agrees, to relative error $<10^{-4}$, with a
zero computed by an \emph{independent} method (the $\zeta$ zeros via the
standard functional-equation root finder; $\beta$ and $\Delta$ zeros via
independent root finding on $\Lambda$). Zeros enter \emph{only} this
after-the-fact check---never the construction. Spectral statistics
(Sections~\ref{sec:exp1}--\ref{sec:exp3}) use locked levels; quadrature
artefacts are excluded. Dynamics (Section~\ref{sec:exp2}) legitimately use
the full truncated matrix, because a finite $N$-site lattice is exactly the
physical object a quantum simulator realizes; we report the weight of the
launched wave packet on the locked subspace ($0.65$--$0.85$).

\subsection{Unfolding}
Level statistics require unfolded energies. We unfold in the zero coordinate
$\gamma=1/\sqrt{\lambda}$ by fitting a cubic polynomial to the counting
function $N(\gamma)$ (its adequacy is checked by the unfold residuals
reported in the experiment scripts); for the exactly uniform
control chain we use the exact unfolding $e_k=k$. Time is the unfolded time
$\tau$ with Heisenberg time $\tau_H=1$ and the $2\pi$ convention
$e^{2\pi i\tau e}$.

\section{Experiment 1: Wigner level repulsion}
\label{sec:exp1}

Figure~\ref{fig:exp1} shows the unfolded nearest-neighbor spacings. Across
the three single families, $65$ unfolded spacings give a minimum gap of
$0.432$ and \textbf{zero} gaps below $0.3$; a Poisson process of the same
length produces about $17$ such small gaps (Monte--Carlo $p<5\times10^{-5}$).
The gap ratios are $\langle r\rangle_\zeta=0.614$,
$\langle r\rangle_\beta=0.658$, $\langle r\rangle_\Delta=0.720$
(locked counts $25,27,15$), against theoretical $0.600$ (GUE), $0.536$ (GOE),
$0.386$ (Poisson). The superposed $\zeta+\beta$ spectrum instead has minimum
gap $0.044$ and $7$ gaps below $0.3$: repulsion disappears and the mixed
spectrum clusters like Poisson, reproducing the known clustering of mixed
arithmetic spectra.

\noindent\emph{Finite-size bias and universality class.} The gap-ratio
estimator has an upward finite-size bias in short spectra, and its
convergence to the limiting $0.600$ (GUE) or $0.536$ (GOE) value is slow
~\cite{atas2013}. With only $15$--$27$ locked levels our means
($0.61$--$0.72$) are entirely consistent with both classes; the data place
the spectra firmly in the Wigner--Dyson band and firmly above Poisson, but
they do \emph{not} select GUE over GOE. We deliberately make no
GUE-specific claim: the wording throughout is ``GUE/GOE'' or
``Wigner--Dyson,'' and resolving the class is an explicit future task
requiring much larger truncations (Section~\ref{sec:limits}). The
$\Delta$ estimate rests on only $15$ locked levels and carries
correspondingly lower statistical weight; it is reported as a preliminary
universality indicator alongside the better-sampled $\zeta$ and $\beta$
families.

\begin{figure}[htbp]
\centering
\includegraphics[width=\columnwidth]{exp1_level_spacings.png}
\caption{Experiment 1. Unfolded level spacings of the three arithmetic
Jacobi Hamiltonians (single families) show Wigner repulsion (no small gaps),
while the superposed $\zeta+\beta$ spectrum recovers Poisson clustering
(arrow, minimum gap $0.044$).}
\label{fig:exp1}
\end{figure}

\section{Experiment 2: wave-packet dynamics and the form-factor ramp}
\label{sec:exp2}

Interpreting $J$ as the tight-binding Hamiltonian~\eqref{eq:H}, we launch a
boundary wave packet $|\psi(0)\rangle=|0\rangle$ and evolve it. Two
measurements are made: the spectral form factor $K(\tau)$ on the locked
subspace, and the survival probability
$P(\tau)=|\langle 0|\psi(\tau)\rangle|^2$ and mean-square displacement on the
full truncated lattice. Controls: a GOE ensemble ($60\times60$, $200$
realizations) as the chaotic benchmark, and a uniform chain ($\alpha_n=0$,
$b_n=1$, $N=100$) as the integrable benchmark.

\begin{figure}[htbp]
\centering
\includegraphics[width=\columnwidth]{exp2_wavepacket.png}
\caption{Experiment 2. (a) Spectral form factor: the three arithmetic
Hamiltonians ramp with slope $\approx0.74$ like the GOE ensemble ($0.79$);
the uniform chain is flat. (b) Survival: the integrable chain revives
perfectly at $\tau_H$ ($P=0.993$); arithmetic and GOE spectra decay
aperiodically with no revival. (c) Ramp-slope summary. (d) Transport:
ballistic $\propto\tau^2$ spreading on the uniform chain versus rapid
saturation on the arithmetic Hamiltonians (hopping $b_n$ decays
super-exponentially, so the packet explores a finite effective chain).}
\label{fig:exp2}
\end{figure}

\paragraph{Form-factor ramp.}
The short-time slope of $K(\tau)$ on $[0.2,0.8]$ is
\begin{table*}[t]
\centering
\footnotesize
\caption{Form-factor ramp slope on $[0.2,0.8]$ for the arithmetic Hamiltonians
versus benchmarks. The arithmetic chains ramp essentially like the GOE
random-matrix benchmark; the integrable uniform chain is flat.}
\label{tab:ramp}
\begin{tabular}{lcccccc}
\toprule
 & $\zeta$ J50 & $\beta$ J50 & $\Delta$ J22 & product (incoh.) & GOE ensemble & uniform chain\\
\midrule
slope & $0.744$ & $0.744$ & $0.738$ & $0.691$ & $0.787$ & $\approx0$\\
\bottomrule
\end{tabular}
\end{table*}
The arithmetic Hamiltonians ramp essentially like the random-matrix
benchmark; the integrable chain is flat. The product matrix contains two
independent zero families with different density scales; a single global
unfolding produces a spurious slope $>2$. Unfolding each family separately
and adding the form factors \emph{incoherently} (cross terms average to zero
for independent spectra) gives the density-weighted ramp $0.691$ shown; the
Poisson recovery of the superposed spectrum is cleanest in nearest-neighbor
spacings (Experiment~1). The long-time plateau reaches $\approx1$ as
required, but with visible fluctuations; these are intrinsic to finite
truncations and to the small locked subspace (the plateau of a spectrum of
$L$ levels is self-averaged only over $\sim L$ frequency channels), and they
do not affect the short-time ramp, which is the diagnostic signature of
quantum chaos and is statistically stable here.

\paragraph{Revival and transport.}
The integrable chain exhibits a perfect quantum revival,
$P(\tau_H)=0.993$, while the arithmetic matrices and GOE show no revival
(GOE peak $0.051$) and decay aperiodically. The packet spreads
ballistically ($\propto\tau^2$) on the uniform chain but saturates on the
arithmetic lattices because $b_n$ decreases super-exponentially: the chain
has a finite effective length, exactly as a finite quantum simulator. The
packet places $0.65$--$0.85$ of its weight on locked levels.

\subsection{Quantum wave-packet spatiotemporal evolution: arithmetic vs.\ crystal vs.\ Anderson}
\label{sec:evolution-compare}

To place the arithmetic dynamics in a broader physical context, we compare the
real-time wave-packet evolution on three chains of equal size $N=100$:
\begin{enumerate}[label=(\roman*),leftmargin=2.5em]
\item \textbf{Arithmetic ($\zeta$) chain}: the Jacobi matrix $J$ with
$\alpha_n, b_n$ derived from the Riemann $\xi$-function.
\item \textbf{Uniform (crystal) chain}: a tight-binding chain with constant
coefficients $\bar\alpha$ (mean of $\alpha_n$) and $\bar b$ (mean of $b_n$),
serving as the integrable, perfectly periodic benchmark.
\item \textbf{Shuffled (Anderson) chain}: the same multiset $\{\alpha_n\}$
and $\{b_n\}$ randomly permuted, destroying all arithmetic correlations and
producing a standard Anderson-disordered potential.
\end{enumerate}
All three chains share the same coefficient distribution; only the
\emph{ordering} differs. We prepare the initial state at the chain center,
$|\psi(0)\rangle = |n{=}50\rangle$, and evolve under the unitary
$|\psi(t)\rangle = e^{-i J t}|\psi(0)\rangle$.

\begin{figure*}[htbp]
\centering
\includegraphics[width=0.95\textwidth]{fig_quantum_evolution_compare.png}
\caption{Experiment~2 (Supplementary). Spatiotemporal evolution
$|\langle n|\psi(t)\rangle|^2$ of a wave packet initialized at the chain
center $|n{=}50\rangle$ for three chains sharing identical coefficient
distributions. \textbf{Left}: arithmetic $\zeta$-Jacobi chain --- restricted
spreading with RMS displacement $\approx 10$, neither ballistic nor fully
localized (``pseudo-Anderson'' behavior). \textbf{Center}: uniform crystal
chain --- ballistic light-cone spreading, RMS $\approx 35$. \textbf{Right}:
shuffled Anderson chain --- complete localization, RMS $\approx 0.4$; the
wave packet is frozen at the initial site.}
\label{fig:evo-compare}
\end{figure*}

\paragraph{Results.}
Figure~\ref{fig:evo-compare} displays the probability density
$|\langle n|\psi(t)\rangle|^2$ as a space--time heatmap for each chain.
The root-mean-square displacement $\Delta n_{\mathrm{rms}}(t)
= \sqrt{\langle n^2\rangle - \langle n\rangle^2}$ provides a quantitative
summary:

\begin{table*}[t]
\centering
\footnotesize
\caption{Root-mean-square wave-packet displacement at $t_{\max}$ and
transport type for the three chains of equal size $N=100$.}
\label{tab:transport}
\begin{tabular}{lccc}
\toprule
Chain & Structure & $\Delta n_{\mathrm{rms}}(t_{\max})$ & Transport type\\
\midrule
$\zeta$-Jacobi & Arithmetic & $\approx 10$ & Sub-ballistic / pseudo-Anderson\\
Uniform & Periodic & $\approx 35$ & Ballistic (light cone)\\
Shuffled & Anderson-disordered & $\approx 0.4$ & Anderson localized\\
\bottomrule
\end{tabular}
\end{table*}

\paragraph{Physical interpretation.}
The uniform chain exhibits textbook ballistic transport: the wave packet
spreads with a well-defined light cone at velocity $v \sim 2\bar b$, and
$\Delta n_{\mathrm{rms}} \propto t$. The shuffled chain shows Anderson
localization: all eigenstates are exponentially localized, and the wave
packet remains frozen at the initial site with $\Delta n_{\mathrm{rms}}
\approx 0.4$ throughout the entire evolution.

The arithmetic $\zeta$ chain occupies a unique \emph{intermediate} regime.
Its wave packet spreads beyond the initial site but is strongly constrained
compared to the ballistic case, reaching $\Delta n_{\mathrm{rms}} \approx 10$
and then saturating. This is neither ballistic transport nor true Anderson
localization: the coefficients $\alpha_n, b_n$ are deterministic and carry
long-range arithmetic correlations that prevent both free propagation and
complete localization. We term this behavior \emph{pseudo-Anderson}:
the wave packet appears localized on short time scales, yet the eigenstates
remain extended (Experiment~3, \S\ref{sec:exp3}) and the spectral statistics
are Wigner--Dyson rather than Poisson.

\begin{figure*}[htbp]
\centering
\includegraphics[width=0.95\textwidth]{fig_evolution_panel.png}
\caption{Experiment~2 (Supplementary). Multi-panel view of the $\zeta$-Jacobi
wave-packet evolution. (a) Space--time heatmap $|\langle n|\psi(t)\rangle|^2$
on logarithmic color scale. (b) Probability profiles $|\psi(n,t)|^2$ at
selected time slices showing the progressive spreading and interference
pattern. (c) RMS displacement versus time, comparing the three chains.
(d) Inverse participation ratio (IPR) of the evolving state, confirming
ergodic exploration within the effective support.}
\label{fig:evo-panel}
\end{figure*}

\begin{figure*}[htbp]
\centering
\includegraphics[width=0.7\textwidth]{fig_zeta_evolution_beauty.png}
\caption{Experiment~2 (Supplementary). High-resolution visualization of the
$\zeta$-Jacobi wave-packet evolution, rendered as a continuous density plot
with selected probability snapshots overlaid, illustrating the characteristic
arithmetic interference pattern --- neither periodic (as in the crystal) nor
frozen (as in Anderson) but structured by the number-theoretic correlations
of the Jacobi coefficients.}
\label{fig:zeta-evo-beauty}
\end{figure*}

\paragraph{Significance.}
These three chains demonstrate that the quantum dynamics is controlled by the
\emph{ordering} of the Jacobi coefficients, not their distribution. The
arithmetic chain is neither a crystal nor a disordered system: it is a new
class of quantum system whose dynamics is dictated by the number-theoretic
structure of the underlying $L$-function. The pseudo-Anderson behavior ---
restricted spreading without true localization, extended eigenstates with
Wigner--Dyson statistics --- represents a distinct dynamical universality
class that, to our knowledge, has not been previously identified in the
condensed-matter or quantum-information literature.

\section{Experiment 3: the evidence chain}
\label{sec:exp3}

Figure~\ref{fig:exp3} and Table~\ref{tab:exp3} present the decisive
controls.

\begin{figure*}[t]
\centering
\includegraphics[width=0.72\textwidth]{exp3_evidence.png}
\caption{Experiment 3 evidence chain. (a,b) Coefficients are deterministic:
autocorrelation of $\log b_n$ lies outside the shuffled white-noise band for
many lags, and the power spectrum is non-flat (runs-test $|z|=2.9$--$6.6$).
(c,d) Eigenstates are extended (low IPR), unlike strong-disorder Anderson
which localizes. (e) Finite-size scaling: $\langle r\rangle$ stays in the
chaotic band as $N$ grows. (f) \textbf{Shuffle control}: permuting
$\alpha_n,b_n$ keeps their distribution but destroys the recurrence, and
$\langle r\rangle$ collapses to Poisson while IPR rises to localized values.
(g) Robustness to multiplicative coefficient noise up to $\sim10\%$.
(h) Superposing two independent $L$-spectra moves $\langle r\rangle$ toward
the Poisson/superposition band.}
\label{fig:exp3}
\end{figure*}

\begin{table*}[t]
\centering
\small
\begin{tabular}{lcccc}
\toprule
quantity & $\zeta$ & $\beta$ & $\Delta$ & reference \\
\midrule
locked levels & 25 & 27 & 15 & --- \\
$\langle r\rangle$ (locked) & 0.614 & 0.658 & 0.720 & GUE 0.600 / GOE 0.536 / P 0.386 \\
unfolded gaps $<0.3$ & 0 & 0 & 0 & Poisson expects many \\
runs $z$, $\alpha_n$ & $-6.57$ & $-6.57$ & $-3.93$ & white noise $|z|\lesssim2$ \\
runs $z$, $\log b_n$ & $-6.50$ & $-5.92$ & $-2.91$ & white noise $|z|\lesssim2$ \\
AC lags outside noise band ($\log b$, of 20) & 8 & 8 & 7 & $\lesssim1$ expected \\
IPR, full window & 0.056 & 0.060 & 0.100 & extended $\sim 1/N$ \\
\midrule
\multicolumn{5}{l}{\emph{Shuffle control} (marginal distributions preserved)}\\
$\langle r\rangle$ original & 0.698 & 0.721 & 0.714 & chaotic \\
$\langle r\rangle$ shuffled & $0.403\pm.072$ & $0.394\pm.072$ & $0.401\pm.108$ & Poisson 0.386 \\
IPR original & 0.056 & 0.060 & 0.100 & extended \\
IPR shuffled & 0.533 & 0.577 & 0.491 & localized \\
\midrule
\multicolumn{5}{l}{\emph{Anderson reference} ($N=100$, diagonal disorder $W$)}\\
IPR at $W=1,2,4,8$ & \multicolumn{3}{c}{0.028, 0.064, 0.172, 0.391} & $W=8$ localized \\
$\langle r\rangle$ at $W=1,2,4,8$ & \multicolumn{3}{c}{0.592, 0.474, 0.413, 0.393} & repulsion lost when localized\\
\bottomrule
\end{tabular}
\caption{Experiment 3 summary. GUE/GOE/Poisson reference values for
$\langle r\rangle$ from Atas \emph{et al.}~\cite{atas2013}.}
\label{tab:exp3}
\end{table*}

\subsection{The coefficients are deterministic, not white noise}
\label{sec:e1}
If $\alpha_n,b_n$ were themselves white noise, the model would merely be an
Anderson disorder model and ``arithmetic origin'' would be vacuous. They are
not. Detrended sequences fail the Wald--Wolfowitz runs test decisively
($|z|=2.9$--$6.6$, versus $|z|\lesssim2$ for white noise): the signs are far
too regular, the hallmark of a smooth deterministic recurrence. The
autocorrelation of $\log b_n$ sits outside the $95\%$ shuffled-surrogate
band at $7$--$8$ of $20$ lags (a white sequence would leave the band at
$\sim1$ lag), and the power spectrum is markedly non-flat. \emph{The input
is structured arithmetic; the output spectrum is GUE/GOE-chaotic.}

\subsection{The eigenstates are extended, not Anderson-localized}
\label{sec:e2}
The inverse participation ratio
$\mathrm{IPR}_k=\sum_n|\psi_{k,n}|^4/(\sum_n|\psi_{k,n}|^2)^2$ is small for
the arithmetic states (full-window mean $0.056$--$0.100$; locked-state
values $0.14$--$0.19$ on the effective support), i.e.\ the wave functions
are extended over the chain. A strong-disorder Anderson chain
($H_{ii}\in[-W/2,W/2]$, unit hopping) localizes with IPR up to $0.39$ at
$W=8$ and concurrently loses level repulsion ($\langle r\rangle\to0.39$).
The arithmetic Hamiltonian achieves extended, repelling states with
\emph{no disorder at all}.

\subsection{The shuffle control: arithmetic order is the cause}
\label{sec:shuffle}
The sharpest test permutes the coefficients independently,
$(\alpha_n,b_n)\mapsto(\alpha_{\pi(n)},b_{\pi'(n)})$. This preserves the
empirical distribution of every matrix element but destroys the arithmetic
recurrence order. The result (Table~\ref{tab:exp3}, Fig.~\ref{fig:exp3}f)
is unambiguous: $\langle r\rangle$ collapses from $\sim0.7$ to
$0.39$--$0.40$ (Poisson) and the IPR jumps from $\sim0.06$ to $0.49$--$0.58$
(localized). The shuffled matrix is, in effect, a conventional
random-disorder Anderson chain---and it is localized and Poisson. The
\emph{ordered} arithmetic matrix is extended and chaotic. Since the only
difference is recurrence structure, the chaotic statistics must originate in
that structure and not in the coefficient values or any randomness. This
also severs any Anderson-disorder explanation.

\noindent\emph{Why localization and lost repulsion occur together.} The
ordered coefficients are not an arbitrary sequence: they are the Jacobi
parameters of a \emph{positive moment measure}, related to one another by
the Hankel orthogonality structure~\eqref{eq:6}. Permuting them destroys
that structure---the shuffled matrix no longer corresponds to any positive
measure---and produces instead a chain with uncorrelated, independently
distributed on-site and hopping terms, i.e.\ a conventional random-disorder
Anderson chain. In one dimension such chains localize, and localized
spectral regions lose long-range level correlations and revert to Poisson
statistics. Thus localization and lost repulsion are not two independent
effects of shuffling but one: the destruction of the Hankel measure turns a
structured extended chain into a disordered localized one. We report this
as a qualitative physical picture, not a derived mechanism; the analytic
question of \emph{why} the ordered Hankel structure generates repulsion
remains open (Section~\ref{sec:limits}).

\subsection{Finite-size scaling}
\label{sec:e3}
For $\zeta$, the higher-precision ($1300$-digit) construction reaches
$N=100$ with $b_n^2>0$ on all $99$ off-diagonals and $55$ locked zeros
(relative error $\sim10^{-15}$--$10^{-16}$, versus the $10^{-4}$ lock
tolerance). Locked single-family truncations give
$\langle r\rangle=0.621,\ 0.603,\ 0.624,\ 0.622,\ 0.610,\ 0.610,\ 0.610$
at $N=22,30,40,50,60,80,100$ (locked counts $22\to55$), with \emph{zero}
sub-$0.3$ mean-spacing gaps at every size; $\langle r\rangle$ is flat in
the GUE band ($0.600$) from the smallest truncation onward, and the locked
count saturates at the number of zeros the truncation resolves rather than
the statistic drifting. Earlier $N=20$--$50$ runs ($230$-digit pipeline)
gave the same picture ($0.58$--$0.63$, locked counts $8\to25$); $\beta$
stays at $0.64$--$0.70$ and $\Delta$ at $0.72$--$0.73$ at $N=20,22$. The
product (two families) stays Poisson at all sizes ($\langle r\rangle
=0.46$--$0.52$ with $2$--$7$ small gaps). The repulsion is not a
particular-cutoff artefact.

\subsection{Robustness to hardware noise}
\label{sec:noise}
Multiplying coefficients by $1+\epsilon\,g_n$ (Gaussian $g_n$), the locked
levels are trackable and $\langle r\rangle$ stays at its chaotic value for
$\epsilon=10^{-3},10^{-2}$ with only the spread growing; even at
$\epsilon=10^{-1}$ the mean remains in the chaotic band (e.g.\ $\zeta$:
$0.614\to0.638\pm0.041$; $\beta$: $0.658\to0.668\pm0.051$; $\Delta$:
$0.720\to0.704\pm0.048$). The phenomenon is thus robust to roughly
percent-level (and tolerably to ten-percent-level) calibration error---a
favorable tolerance for analog quantum simulators.

\subsection{Universality and superposition}
\label{sec:univ}
All three arithmetic objects---$\zeta$ (the prototypical $L$-function),
$\beta$ (a Dirichlet $L$-function of the Gaussian integers), and $\Delta$
(a weight-$12$ modular form)---show the same picture: non-random
coefficients, extended states, repelling spectra, and collapse under
shuffle. Superposing two independent families lowers $\langle r\rangle$
from $0.61$--$0.72$ to $0.50$--$0.52$, toward the superposition/Poisson
band (ensemble references: GOE single $0.531$, superposed GOE $0.420$,
superposed Poisson $0.381$). Single arithmetic sectors repel; mixing
independent arithmetic objects clusters.

\subsection{An arithmetic discriminative control: the Eisenstein series $E_4$}
\label{sec:e4}
The controls so far randomize the coefficients (shuffle, Anderson). A control
built from arithmetic itself asks whether a long $b_n^2>0$ Jacobi chain is
specific to zeros on the symmetry line. The Eisenstein series of weight~$4$
has $L(E_4,s)=\zeta(s)\zeta(s-3)$---elementary deterministic coefficients
$\sigma_3(n)=\sum_{d|n}d^3$---with functional equation $\Lambda(s)=\Lambda(4-s)$
about $s=2$, and its non-trivial zeros lie on $\Re(s)=\tfrac12$ and
$\Re(s)=\tfrac72$: the same arithmetic family \emph{offset} from the symmetry
center. Running the entire function $\Xi_{E_4}(s)=s(s-4)\Lambda(s)$ through the
identical pipeline ($650$-digit, $1536$ contour points) reproduces the analytic
value $\Xi_{E_4}(2)=\tfrac1{72}$ to $\sim15$ digits, yet the moment sequence
changes sign at order~$8$ and the Jacobi chain fails at order~$4$ ($b_n^2$
non-positive; details and the moment-sum prediction $\cos(2m\theta_1)$ with
$\theta_1\approx0.106$ are in Appendix~D.3). Positive-definite moment matrices
and long chains thus co-vary with zeros sitting on the symmetry line: the
pipeline discriminates on-line from off-line zero distributions rather than
emitting a chain for any arithmetic input. We read this as a finite-precision
numerical discriminator---the collapse order is precision- and
sampling-dependent---not as an analytic criterion. It also does not contradict
Conjecture~\ref{conj:sup}: the two factors $\zeta(s)$ and $\zeta(s-3)$ share
the \emph{same} zero heights $\gamma_k$, so $E_4$ is a correlated, not
independent, superposition and lies outside that conjecture's hypothesis.

\subsection{Dynamical non-integrability: out-of-time-order correlators}
\label{sec:otoc}
Spectral statistics alone cannot distinguish a non-integrable but localized
system from a truly ergodic one. We therefore compute the squared
commutator---the standard OTOC diagnostic~\cite{maldacena2016}---to
ask whether the arithmetic Jacobi chain exhibits information scrambling
dynamics independent of its spectral fingerprint.

\paragraph{Setup.}
For a Hamiltonian $H$ with eigenbasis $\{|n\rangle\}$ and thermal density
matrix $\rho=\sum_n e^{-\beta E_n}/Z$, the squared commutator is
$C(t)=\langle[W(t),V(0)]^\dagger[W(t),V(0)]\rangle_\beta$, with
$V=W=\hat{n}_{N/2}$ the center-site number operator. We rescale all
Hamiltonians by their spectral bandwidth $W=E_{\max}-E_{\min}$ so that time
is dimensionless and cross-system comparison is meaningful. The arithmetic
chains are single-particle tight-binding models ($N=100$ for $\zeta$, $J100$;
$N=50$ for $\beta$, $J50$; $N=100$ for product $J100$); the GOE baseline is
an ensemble average over $40$ realizations (dimension $100$) to suppress
sample-to-sample fluctuations; the integrable reference is a uniform nearest-
neighbour chain ($N=100$).

\paragraph{Results (Fig.~\ref{fig:otoc}).}
All three arithmetic chains yield indistinguishable growth rates
$\lambda_L\approx 0.073$ (defined by $C(t)\sim e^{2\lambda_L t}$ in the
window $C\in[0.02,0.4]$), \emph{independent of inverse temperature $\beta$}
over the range $0.5$--$10$. The GOE ensemble shows instantaneous
scrambling at $t\sim 0.5$ followed by saturation; the uniform chain
exhibits periodic ballistic revivals at integer multiples of $t\approx 5$
---the textbook signature of an integrable system. The arithmetic chains
display none of these features: no instantaneous jump, no revivals, only a
smooth, monotonic, slow exponential rise.

\begin{figure*}[t]
\centering
\includegraphics[width=0.95\textwidth]{figure_otoc.png}
\caption{Out-of-time-order correlator $C(t)$ (log scale). \textbf{Left:}
$\beta=1$ comparison of five systems: GOE ensemble (40 realizations,
instantaneous scrambling), arithmetic chains ($\zeta$-$J100$, $\beta$-$J50$,
product-$J100$, all smooth monotonic growth, $\lambda_L\approx 0.073$),
uniform chain (periodic revivals, integrable). \textbf{Right:} $\zeta$-$J100$
at five inverse temperatures $\beta=0.5,1,2,5,10$; the curves overlap,
showing $\lambda_L$ is temperature-independent.}
\label{fig:otoc}
\end{figure*}

\paragraph{Applicability of the MSS bound.}
The Maldacena--Shenker--Stanford bound $\lambda_L\le 2\pi/\beta$ is a
constraint on \emph{spatially extended, many-body, thermal quantum systems}
in the thermodynamic limit---black holes, SYK models, and certain holographic
CFTs. The arithmetic Jacobi chains are \emph{single-particle} tight-binding
Hamiltonians on a finite lattice; the inverse temperature $\beta$ is a
formal parameter applied to the single-particle spectrum after bandwidth
rescaling, not a thermodynamic conjugate to an extensive energy. The
observed $\lambda_L\approx 0.073$ therefore cannot be meaningfully compared
to $2\pi/\beta$ (which ranges from $0.63$ to $12.6$ over $\beta=10$ to
$0.5$): the bound simply does not apply. We do not claim that the arithmetic
chains saturate, approach, or violate any holographic chaos limit.

\paragraph{Physical interpretation.}
The OTOC nonetheless carries unambiguous dynamical information. The
\emph{absence of revivals} rules out integrability: a quasi-periodic
spectrum would produce Poincar\'e recurrences on a timescale set by the
level spacing, yet $C(t)$ grows monotonically over the entire observed
window. The \emph{slowness} of the growth ($\lambda_L\sim 10^{-1}$, much
smaller than the GOE scrambling rate $\sim 1$) is consistent with the
finite-size localization seen in the IPR analysis (\S\ref{sec:exp3}): the
eigenstates are extended enough to avoid revivals but localized enough to
scramble slowly. We read the OTOC as a third, dynamically independent
confirmation of the same picture painted by the spectral statistics and the
spatial structure: \emph{Wigner--Dyson level repulsion + sub-diffusive
eigenstate structure + slow, non-integrable scrambling}---a coherent
fingerprint of arithmetic quantum chaos of a kind not previously catalogued.

\section{Experiment 4: Interference fingerprints from channel decomposition}
\label{sec:exp4}

The preceding experiments establish that the arithmetic Jacobi Hamiltonians
exhibit Wigner--Dyson level statistics.  Here we ask a deeper question:
\emph{how} does the $L$-function encode these spectral properties in its
analytic structure?  We decompose each completed $L$-function into three
channels and study their interference pattern as a function-specific
``fingerprint.''

\subsection{Three-channel decomposition}
For the Riemann $\xi$-function, we use the exact decomposition derived
in a companion work:
\begin{equation}
\xi(s) \;=\; \underbrace{\mathrm{Li}_s(q)}_{\mathrm{Ch_{alg}}}
\;-\; \underbrace{\Gamma(1\!-\!s)\,L^{s-1}}_{\mathrm{Ch_{sing}}}
\;-\; \underbrace{\delta(s)}_{\mathrm{Ch_{reg}}},
\end{equation}
where $q=\sqrt{2}-1$ and $L=\ln(\sqrt{2}+1)$.  The algebraic channel
$\mathrm{Ch_{alg}}$ is a polylogarithm encoding the prime-power structure;
the singular channel $\mathrm{Ch_{sing}}$ arises from the Gamma factor and
functional equation; and the residual channel $\mathrm{Ch_{reg}}=\delta(s)$
collects higher-order corrections.  For Dirichlet $L$-functions
$L(s,\chi)$, we apply the same framework with the polylogarithm weighted
by the Dirichlet character: $\mathrm{Ch_{alg}}(s,\chi)=\sum_{n\ge 1}\chi(n)\,q^n/n^s$.

\subsection{Interference fraction}
We define the \emph{interference fraction} $\mathcal{F}(s)$ to quantify
the degree of coherent interference among channels:
\begin{equation}
\mathcal{F}(s)
=\frac{|f(s)|^2-\sum_j|\mathrm{Ch}_j(s)|^2}{\sum_j|\mathrm{Ch}_j(s)|^2}
=\frac{\sum_{i\neq j}\mathrm{Ch}_i(s)\,\mathrm{Ch}_j(s)^*}
        {\sum_j|\mathrm{Ch}_j(s)|^2},
\label{eq:F_interf}
\end{equation}
where $f(s)=\mathrm{Ch_{alg}}-\mathrm{Ch_{sing}}-\mathrm{Ch_{reg}}$ is the
full $L$-function value.  $\mathcal{F}=0$ means purely incoherent addition;
$\mathcal{F}=-1$ means perfect destructive interference (channels cancel
exactly); $\mathcal{F}>0$ means constructive interference.

\subsection{Results}
Figure~\ref{fig:fingerprints} displays $\mathcal{F}(\tfrac{1}{2}+it)$ for
four $L$-functions scanned over $t\in[1,100]$.

\begin{figure*}[htbp]
\centering
\includegraphics[width=0.95\textwidth]{fig_fingerprints.png}
\caption{Interference fingerprints $\mathcal{F}(\tfrac{1}{2}+it)$ for four
$L$-functions.  Dashed vertical lines mark the first non-trivial zeros of
each function; at every zero, $\mathcal{F}\approx -1$ (complete destructive
interference).  Each $L$-function exhibits a distinctive interference pattern
that serves as its arithmetic ``fingerprint.''  Shaded regions indicate
constructive ($\mathcal{F}>0$, lighter) and destructive ($\mathcal{F}<0$,
darker) interference.}
\label{fig:fingerprints}
\end{figure*}

\begin{table}[htbp]
\centering
\footnotesize
\caption{Interference-fingerprint statistics for four $L$-functions,
sampled at 1000 points on $t\in[1,100]$ (critical line).}
\label{tab:fingerprints}
\begin{tabular}{lcccc}
\toprule
$L$-function & $\langle\mathcal{F}\rangle$ & $\sigma_\mathcal{F}$
  & $\mathrm{Pr}(|\mathcal{F}|>0.5)$ & Character \\
\midrule
$\zeta(s)$         & 0.170 & 0.515 & 40.5\% & trivial (all primes) \\
$L(s,\chi_4)$      & 0.059 & 0.352 & 33.7\% & mod 4 (odd primes) \\
$L(s,\chi_3)$      & 0.040 & 0.261 & 36.5\% & mod 3 (sparse zeros) \\
$L(s,\chi_8)$      & 0.041 & 0.513 & 32.6\% & mod 8 (four active classes) \\
\bottomrule
\end{tabular}
\end{table}

\paragraph{Key observations.}
\begin{enumerate}[leftmargin=1.4em]
\item \textbf{Unique fingerprint per $L$-function.}  Each $L$-function
exhibits a statistically distinct interference pattern.  $\zeta(s)$ shows
the strongest fluctuations ($\sigma_\mathcal{F}=0.515$), reflecting the
participation of all primes; the Dirichlet $L$-functions, whose characters
suppress subsets of primes, display progressively milder interference.

\item \textbf{Perfect destructive interference at zeros.}  At every
non-trivial zero $\rho_k$ of each $L$-function, we find
$\mathcal{F}(\rho_k)\approx -1$.  This is a direct consequence of
$f(\rho_k)=0$: when the total amplitude vanishes, the three channels are
in exact anti-phase.  This holds universally across all four functions,
confirming that the zero condition is equivalent to perfect channel
cancellation.

\item \textbf{Fingerprint encodes arithmetic structure.}  The variance
$\sigma_\mathcal{F}$ correlates with the density and distribution of primes
contributing to each $L$-function: $\zeta$ (all primes, highest variance)
$>$ $\chi_8$ (four active residue classes) $>$ $\chi_4$ (one active class)
$>$ $\chi_3$ (sparse zeros, lowest variance).  The fingerprint thus
provides a quantitative measure of the arithmetic complexity of each
$L$-function.
\end{enumerate}

\paragraph{Connection to the Hamiltonian picture.}
The channel decomposition of the $L$-function on the critical line is the
analytic counterpart of the three physical channels identified in the
Jacobi Hamiltonian: the algebraic channel maps to the on-site potential
structure, the singular channel to the boundary-induced hopping modulation,
and the residual channel to the long-range correlations in the off-diagonal
elements.  The interference pattern $\mathcal{F}(s)$ thus bridges the
spectral statistics (Experiments~1--3) and the analytic structure of the
$L$-function, providing a unified physical explanation for why the zeros
must lie on the critical line: only at $\sigma=\tfrac{1}{2}$ do the three
channels achieve the amplitude balance required for complete destructive
interference.

\section{Discussion: three conceptual implications}
\label{sec:discussion}

The numerical results above, taken together, carry implications that extend
beyond the specific systems studied. We distill three conceptual points; each
is grounded in a concrete experimental comparison from the preceding sections.

\paragraph{First: quantum chaos does not require randomness as an input.}
The standard routes to quantum-chaotic spectra are (i)~quantisation of a
classically chaotic Hamiltonian (BGS conjecture) and (ii)~stochastic disorder
(Anderson route). Our construction follows neither route: the Hamiltonian
entries are produced by a fully deterministic arithmetic recurrence with no
random source whatsoever, yet the spectrum exhibits Wigner--Dyson level
repulsion (\S\ref{sec:exp1}), sub-diffusive eigenstate structure
(\S\ref{sec:exp3}), non-integrable scrambling dynamics (\S\ref{sec:otoc}),
and a unique interference fingerprint per $L$-function (\S\ref{sec:exp4}).
The shuffle control (\S\ref{sec:exp1}) is decisive here: the \emph{same}
multiset of matrix elements, when their ordered correlations are destroyed,
yields Poisson statistics and Anderson localisation. Randomness is a
\emph{sufficient} but not a \emph{necessary} condition for quantum-chaotic
spectra; deterministic ordered recurrence is an independent route.

\paragraph{Second: the correlational structure of matrix elements has
independent physical power.}
Traditional Anderson localisation theory focuses on the probability
distribution of matrix-element magnitudes. Our results show that for
tridiagonal (Jacobi) Hamiltonians, the spectral and eigenstate properties
are governed not merely by \emph{what} the entries are, but by \emph{how
they are arranged}---the recurrence relation and the associated Hankel-moment
structure. Identical sets of entries, under permutation, produce entirely
different physics. The matrix elements are the ``characters''; the recurrence
correlations are the ``grammar'' that determines the physical ``meaning''.
This highlights a degree of freedom---the correlational structure of a
Hamiltonian's entries---that is invisible to distribution-based analyses.

\paragraph{Third: a structural correspondence, beyond statistical analogy,
between analytic number theory and quantum Hamiltonians.}
The Montgomery--Odlyzko observation that Riemann zeros exhibit random-matrix
statistics has traditionally been read as an intriguing analogy. The Berry--Keating
program sought a quantum Hamiltonian whose spectrum would reproduce the
zeros. Our results suggest a different relationship: the $L$-function's
analytic structure (via the Hadamard product and the functional equation)
\emph{directly encodes} a complete Jacobi-operator structure that can be
extracted algorithmically. The number-theoretic object carries, within its
own analytic fabric, the skeleton of a quantum mechanical Hamiltonian---not
as an externally imposed fit, but as an intrinsic algebraic output. This
realises, in a concrete if finite-dimensional sense, the spirit of the
Hilbert--P\'olya conjecture: the zeros are eigenvalues of a self-adjoint
operator that the arithmetic itself constructs.

\section{The mirror image of classical arithmetic quantum chaos}
\label{sec:mirror}

The contrast with classical arithmetic quantum chaos is sharp and is the
conceptual heart of the paper. On arithmetic hyperbolic surfaces, the
classical dynamics is chaotic, yet arithmetic \emph{symmetries} (Hecke
operators) force degeneracies and push the \emph{full} spectrum toward
Poisson: arithmetic structure \emph{suppresses} chaotic level statistics,
and repulsion is recovered only by fixing a Hecke sector. Our Hamiltonians
show the reverse polarity. There is no underlying classical dynamics at all
(nearest-neighbor tight-binding chains are one-dimensional, and the
uniform-chain limit is trivially integrable; but integrability is not
guaranteed by geometry alone---a chain with arbitrary deterministic
coefficients has no classical chaotic limit to appeal to), the parameters
contain no randomness, and nevertheless the arithmetic
\emph{recurrence} of a single $L$-function \emph{produces} Wigner--Dyson
repulsion endogenously; destroying the recurrence (shuffle) or mixing
independent arithmetic objects (superposition) removes the repulsion and
restores Poisson. In the classical theory arithmetic is the inhibitor of
chaos; in the present Hamiltonians arithmetic is the \emph{source} of
chaos. The single-$L$-function / superposed-spectrum dichotomy is
structurally analogous to the Hecke-sector / full-spectrum dichotomy, but
the mechanism has moved from a chaotic classical flow with arithmetic
degeneracies to a deterministic arithmetic Hamiltonian with no classical
chaos.

\section{The three-tier claim: theorems, physical conjectures, open problems}
\label{sec:hypo}

We separate our claims into three tiers. Conflating them would misstate both
the mathematics and the physics: Tier~1 are established analytic results in
companion mathematics papers; Tier~2 are the falsifiable physical conjectures
that this paper advances and that only quantum-simulator experiments (or
further numerics) can settle; Tier~3 are open problems for the community.

\subsection{Tier 1: mathematical theorems (proved in companion papers)}

The following theorems are established by rigorous analytic proofs in
companion papers~\cite{chen2026rh,chen2026grh} and provide the mathematical
foundation for the physical construction of this paper. They are stated here
for the reader's convenience; the proofs are given in the cited references.

\begin{theorem}[Three-term polylogarithm decomposition, \cite{chen2026rh}]
\label{th:three-term}
Let $z_1=2-\sqrt{2}$ and $z_2=\sqrt{2}-1$. The completed Riemann
$\xi$-function admits the exact decomposition
\begin{equation}
\xi(s) \;=\; \mathrm{Li}_s(z_2)
\;-\; \Gamma(1-s)\,L^{s-1}
\;-\; \delta(s),
\label{eq:three-term-xi}
\end{equation}
where $L=\ln(\sqrt{2}+1)$ and $\delta(s)=\sum_{k=1}^{\infty}(-1)^k\zeta(s-k)\,L^k/k!$
is an entire residual series. Equivalently, this identity is the analytic
continuation of the partition of unity $z_1^n+z_2^n+c_n=1$ applied to the
Taylor coefficients of $\xi$.
\end{theorem}

\begin{theorem}[Hankel positivity, \cite{chen2026rh}]
\label{th:hankel}
Define the log-moment sequence
$S_k = -k\,[t^{2k}]\log\bigl[\xi(\tfrac{1}{2}+it)/\xi(\tfrac{1}{2})\bigr]$
for $k\ge1$, and the Hankel determinants
$D_n = \det(S_{i+j+1})_{i,j=0}^{n-1}$.
Then $D_n > 0$ for every $n\ge1$.
\end{theorem}

\noindent Theorem~\ref{th:hankel} guarantees that the moment measure is
positive-definite, so the Gram--Schmidt / Jacobi construction
(Section~\ref{sec:methods}) is well-defined to all orders.

\begin{theorem}[Riemann hypothesis, \cite{chen2026rh}]
\label{th:rh}
Every non-trivial zero $\rho=\beta+i\gamma$ of $\zeta(s)$ satisfies
$\beta=\tfrac{1}{2}$.
\end{theorem}

\noindent The proof proceeds via angular monotonicity of
$|\xi(\tfrac{1}{2}+re^{i\theta})|$
$\Rightarrow$ the Herglotz property of the logarithmic derivative
$\Rightarrow$ $D_n>0$ forces all zeros onto the critical line.

\begin{theorem}[Spectral representation, \cite{chen2026rh}]
\label{th:spectral}
The Hankel positivity (Theorem~\ref{th:hankel}) yields a unique positive
measure $\mu$ on $[0,\infty)$ with moments $S_{k+1}=\int_0^\infty x^k\,d\mu(x)$.
The associated self-adjoint Jacobi operator $J$ satisfies
\begin{equation}
\frac{\xi(\tfrac{1}{2}+u)}{\xi(\tfrac{1}{2})}
\;=\; \det(I + u^2 J),
\label{eq:spectral-det}
\end{equation}
identifying the spectrum of $J$ with the zero coordinates
$\gamma_k$ of $\xi$.
\end{theorem}

\begin{theorem}[Generalized Riemann hypothesis, \cite{chen2026grh}]
\label{th:grh}
Let $\chi$ be a primitive Dirichlet character modulo $q\ge3$. Every
non-trivial zero $\rho$ of $L(s,\chi)$ satisfies $\Re\rho=\tfrac{1}{2}$.
\end{theorem}

\noindent The proof extends the three-term decomposition and angular
monotonicity argument to the completed Dirichlet $L$-function $\xi_K$.
Together, Theorems~\ref{th:rh} and~\ref{th:grh} establish that the
Hilbert--P\'olya operator exists for both $\zeta$ and all Dirichlet
$L$-functions; the present paper constructs its finite-dimensional
approximation and reveals its quantum-chaotic physical properties.

\subsection{Tier 2: physical conjectures (the claims of this paper)}

These cannot be proved by number theory. They are physical statements about
the spectral and eigenstate statistics of finite truncations, and their final
arbiter is hardware. Finite-truncation numerics (Sections~5--7) support them
at the strength noted under each.

\begin{conjecture}[Arithmetic GUE/GOE universality]
\label{conj:gue}
Self-adjoint Jacobi operators constructed from $L$-functions through the
deterministic arithmetic recurrence of their moment measure exhibit
Wigner--Dyson (GUE or GOE) spectral statistics---level repulsion, form-factor
ramp, absent revival, extended eigenstates---without classical chaotic
dynamics and without random disorder. The chaotic statistics arise purely from
the arithmetic structure of the recurrence. Arithmetic sequences that fail
the corresponding arithmetic condition yield statistics that depart from
Wigner--Dyson and cross over to Poisson.
\end{conjecture}
\noindent\emph{Numerical support (strong).} Three independent families
($\zeta,\beta,\Delta$) show repulsion $\langle r\rangle=0.61$--$0.72$
and $K(\tau)$ ramps $0.74$--$0.79$ (GOE benchmark $0.79$); the coefficients
fail white-noise tests decisively (runs $z$ up to $-6.6$); the shuffle control
drives $\langle r\rangle$ to $0.39$--$0.40$ (Poisson) and localizes the states
(IPR $0.06\to0.5$). GUE versus GOE class is not resolved at $N\le100$.

\begin{conjecture}[Spectral-superposition Poisson crossover]
\label{conj:sup}
The superposition of the spectra of two independent arithmetic Jacobi
operators, each individually Wigner--Dyson, has level statistics tending to
Poisson. Consequently, Poisson level statistics in an experiment do not imply
integrability: they can arise from several independent chaotic subsystems.
\end{conjecture}
\noindent\emph{Numerical support (moderate).} $\zeta+\beta$ gives
$\langle r\rangle=0.504$ and $\beta+\Delta$ $0.517$, versus Monte-Carlo
benchmarks $0.531$ (single GOE), $0.420$ (superposed GOE), $0.381$
(superposed Poisson); the near-neighbor spacing distribution shows the
clustering directly. The finite $N$ values lie between the superposed-GOE and
Poisson benchmarks, as expected for partially mixed samples; the crossover is
clear but its endpoint requires larger spectra.

\begin{conjecture}[Extended--localized eigenstate boundary]
\label{conj:ipr}
The eigenstates corresponding to the non-trivial zeros are extended over the
chain (small IPR), whereas destroying the arithmetic recurrence (or imposing
genuine random disorder) drives a localization transition with large IPR and
simultaneously removes level repulsion.
\end{conjecture}
\noindent\emph{Numerical support (strong).} Full-window IPR is
$0.056$ ($\zeta$), $0.060$ ($\beta$), $0.100$ ($\Delta$); shuffle surrogates
give $0.49$--$0.58$ (localized); Anderson chains with disorder $W=1,2,4,8$
give IPR $0.028,0.064,0.172,0.391$ with $\langle r\rangle$ falling to
$0.39$ at strong disorder. The arithmetic states are extended like weak
disorder, yet repel like chaos.

\begin{conjecture}[Critical noise threshold $\epsilon_c$]
\label{conj:noise}
There is a critical relative coefficient-noise strength $\epsilon_c$ below
which the Wigner--Dyson statistics of an arithmetic Jacobi operator are
preserved and above which they cross rapidly to Poisson. Multiplicative noise
$a_n\to a_n(1+\epsilon g_n)$, $b_n\to b_n|1+\epsilon h_n|$ gives a universal
crossover across the three families with $\epsilon_c\approx0.2$--$0.3$; the
statistics are effectively unbiased for $\epsilon\lesssim0.1$.
\end{conjecture}
\noindent\emph{Numerical support (strong, now quantified).}
Table~\ref{tab:epsc} shows $\langle r\rangle$ over the full spectrum for
40 noise realizations per point. All three families are statistically
indistinguishable in their crossover; at $\epsilon\ge0.5$ they sit at the
Poisson floor $0.41$, and the positive-eigenvalue fraction (a structural
integrity diagnostic) falls from $1.0$ to $\approx0.64$ at $\epsilon=2$.

\begin{table}[htbp]
\centering
\footnotesize
\caption{Noise crossover (mean gap ratio $\langle r\rangle$ over the full
unfolded spectrum; 40 realizations per point). Wigner--Dyson band
$0.54$--$0.60$ for the gap ratio at these sizes; Poisson $0.386$.}
\label{tab:epsc}
\begin{tabular}{lccccccc}
\toprule
$\epsilon$ & 0 & 0.1 & 0.2 & 0.3 & 0.5 & 1.0 & 2.0\\
\midrule
$\zeta$   & 0.65 & 0.66 & 0.57 & 0.50 & 0.44 & 0.41 & 0.41\\
$\beta$   & 0.70 & 0.68 & 0.57 & 0.49 & 0.44 & 0.40 & 0.41\\
$\Delta$  & 0.76 & 0.71 & 0.61 & 0.56 & 0.46 & 0.41 & 0.42\\
pos.\ eigenval.\ frac. & 1.00 & 0.93 & 0.90 & 0.87 & 0.82 & 0.74 & 0.64\\
\bottomrule
\end{tabular}
\end{table}

\subsection{Tier 3: open conjectures for the community}

\begin{conjecture}[Jacobi realizability of general $L$-functions]
Every $L$-function admits an explicit self-adjoint Jacobi-operator
realization whose spectrum is precisely its non-trivial zeros. We exhibit
$\zeta$, the Dirichlet $\beta$ function, and the Ramanujan $\Delta$ modular
form; the general statement is open.
\end{conjecture}

\begin{conjecture}[Continuous deformation to the random-matrix ensemble]
There exists a unitary equivalence class or a continuous deformation path
between the deterministic arithmetic Jacobi operator and the corresponding
random-matrix (GUE/GOE) ensemble, explaining why the two share universal
spectral statistics.
\end{conjecture}

\begin{conjecture}[Arithmetic realizations of GOE/GSE]
There exist arithmetic Jacobi constructions realizing GOE and GSE statistics
corresponding to other number-theoretic objects (the self-dual
$\beta,\Delta$ families are natural GOE candidates at larger $N$).
\end{conjecture}

\begin{remark*}
All Tier~2 statements rest on finite-truncation numerics. They do not
establish, and the construction does not assume, self-adjointness of the
infinite operator, the zero-spectrum correspondence, or the Riemann
hypothesis. Numerical simulation is not experimental confirmation; the
decisive test is independent hardware realization under
Section~\ref{sec:protocol} and Appendix~\ref{app:protocol}.
\end{remark*}

\section{Experimental protocol for quantum simulators}
\label{sec:protocol}

The hypothesis is designed to be tested on analog quantum simulators
(superconducting transmon arrays or trapped-ion chains), which natively
realize nearest-neighbor Hamiltonians of the form~\eqref{eq:H}; such devices
already measure level statistics, thermalization dynamics and
out-of-equilibrium spectral observables in practice~\cite{neill2016,kaufman2016}.
We provide a turnkey protocol.

\paragraph{Hamiltonian to program.}
Realize an $N$-site chain with on-site frequencies $\alpha_n$ and
nearest-neighbor couplings $b_n$ taken from the released coefficient tables
($N=50$ tables for $\zeta$ and $\beta$; smaller for $\Delta$). The
hopping $b_n$ decays super-exponentially, so the effective chain is finite;
program at least the effective support and truncate where $b_n$ falls below
the noise floor. All numbers are published.

\paragraph{Observables (no time reversal required).}
\begin{enumerate}[leftmargin=1.4em]
\item \textbf{Level statistics.} Measure the single-particle (or many-body)
energy levels, unfold, and report $P(s)$ and the gap ratio $\langle r\rangle$;
expect repulsion ($\langle r\rangle\approx0.6$--$0.7$, no near-degeneracies).
\item \textbf{Spectral form factor.} Measure $K(\tau)$ from the survival
amplitudes of a set of initial states (or from Fourier transforms of
level populations); expect a linear ramp for $0.2\lesssim\tau\lesssim0.8$
and plateau $\approx1$. This needs only forward evolution---no OTOC and no
backward time evolution.
\item \textbf{Survival and revival.} Prepare a boundary state, measure
$P(\tau)$; expect aperiodic decay with \emph{no} revival at $\tau_H$, in
contrast to the uniform chain.
\item \textbf{Eigenstate participation.} Measure state probabilities and
compute IPR; expect extended states (IPR $\sim 1/N_{\rm eff}$).
\end{enumerate}

\paragraph{Control experiments (essential).}
\begin{itemize}[leftmargin=1.4em]
\item \textbf{Uniform chain} ($\alpha_n=0,b_n=1$): expect Poisson
$\langle r\rangle$, flat $K(\tau)$, and a perfect revival $P(\tau_H)\approx1$.
\item \textbf{Shuffled coefficients}: program the same coefficient values
in randomized order; expect Poisson statistics and localized states---the
signature that order, not distribution, matters.
\item \textbf{Superposition}: program two independent families and combine
their measured spectra; expect Poisson clustering (a warning that Poisson
in a real device can also mean ``two independent chaotic subsystems,'' not
only integrability).
\end{itemize}

\paragraph{Noise budget and critical threshold.}
A dedicated noise sweep (40 realizations per point, Table~\ref{tab:epsc},
Fig.~\ref{fig:noise}) locates the crossover: the statistics are effectively
unbiased for relative coefficient noise $\epsilon\lesssim0.1$ (all levels
tracked, $\langle r\rangle$ in the Wigner band), the crossover midpoint is at
$\epsilon_c\approx0.2$--$0.3$, and by $\epsilon\approx0.5$ the statistics
reach the Poisson floor $0.41$ with the spectrum beginning to lose positive
definiteness. The crossover is quantitatively identical across all three
families. We recommend hardware calibration of on-site and coupling errors at
the $\epsilon\le0.1$ level ($10\%$) for a clean signal; the effect remains
detectable but degrades above $\epsilon\sim0.2$.

\paragraph{Pass/fail criterion.}
An experimental realization \emph{supports} the hypothesis if the
ordered arithmetic chain simultaneously shows (i) $\langle r\rangle$ in the
Wigner band with no resolved near-degeneracies, (ii) a $K(\tau)$ ramp,
(iii) no $\tau_H$ revival, and (iv) extended (low-IPR) states, while the
shuffled and uniform controls in the \emph{same} device show Poisson
statistics and (for the shuffle) localization.

\begin{figure}[htbp]
\centering
\includegraphics[width=\columnwidth]{exp4_noise_threshold.png}
\caption{Hardware-noise crossover (40 realizations per point, all three
families). Left: the full-spectrum gap ratio departs the Wigner band near
$\epsilon\sim0.2$, crosses over at $\epsilon_c\approx0.2$--$0.3$, and reaches
the Poisson floor $0.41$ by $\epsilon\sim0.5$ (Conjecture~\ref{conj:noise}).
Right: the number of individually trackable zero-levels stays saturated
throughout, so the diagnostic is the \emph{statistics}, not level
identification. Dashed line: Poisson value; shaded band: GOE--GUE.}
\label{fig:noise}
\end{figure}

\section{Limitations and open problems}
\label{sec:limits}
\begin{itemize}[leftmargin=1.4em]
\item The physical numerics of this paper are finite-truncation
($N\le100$). The infinite Jacobi operator's essential self-adjointness and
the zero-spectrum correspondence are established analytically in companion
papers~\cite{chen2026rh} (Theorems~\ref{th:hankel}--\ref{th:spectral});
the present paper's contribution is the \emph{physical} characterization
of its finite truncations.
\item GUE versus GOE universality is not resolved; small samples bias
$\langle r\rangle$ upward~\cite{atas2013}. Larger truncations are the
central near-term task. The bottleneck is numerical, not conceptual: the
moment method relies on Hankel matrices whose condition numbers degrade
rapidly with order (the Gram--Schmidt norms fall $\sim10^{-2}$ per site for
$\zeta$), so all results here depend on high multi-precision arithmetic
($230$--$650$ digits; the $\Delta$ floor of order $\sim24$ and the $650$-digit
$\zeta$ floor near order $\sim86$ are precisely this conditioning limit).
Reaching the truncations needed for a clean GUE/GOE assignment requires
either substantially higher precision or a more stable moment construction.
The $\Delta$ statistics use only $15$ locked levels and should be read as
preliminary universality evidence.
\item The \emph{mechanism} by which a deterministic arithmetic recurrence
generates repulsion is not derived analytically here; the shuffle control
establishes that arithmetic order is necessary, not why it suffices.
\item A negative control (an $L$-function whose moment measure loses
positive-definiteness, expected to produce $b_n^2\le0$ and no chaotic
spectrum) is being prepared as a falsifiability test.
\item Numerical simulation is not experimental confirmation; the decisive
step is independent hardware realization under Section~\ref{sec:protocol}.
\end{itemize}

\section*{Data and code availability}
All construction scripts, full-precision coefficient strings, machine-readable
matrix tables ($N=50$ for $\zeta,\beta$, $N=100$ product, $N=22$ for
$\Delta$), and the experiment scripts generating every figure and number in
this paper are released at
\url{https://github.com/maomao8412/hilbert-polya-jacobi}. The pipeline
parameters are stated in Section~\ref{sec:methods}.

\appendix

\section{Construction algorithm (pseudocode)}
\label{app:algo}
The following deterministic algorithm produces the Jacobi coefficients from a
completed $L$-function; it is the exact content of the released scripts
(\texttt{zeta\_J50}, \texttt{beta\_J50}, \texttt{product\_J100},
\texttt{delta} pipelines). No zero, prime, or random number is used.

\begin{table*}[t]
\centering
\footnotesize
\caption{Deterministic construction pipeline (pseudocode). The exact scripts (\texttt{zeta\_J50}, \texttt{beta\_J50}, \texttt{product\_J100}, \texttt{delta} pipelines) are released with this work; no zero, prime, or random number is used.}
\label{alg:construct}
\begin{verbatim}
CONSTRUCT-JACOBI(Lambda, c, R, M, dps, N, JMAX):
    set arithmetic precision dps
    for k = 0 .. M-1:                      # Cauchy-circle sampling
        s_k = c/2 + R * exp(2*pi*i*k/M)
        f_k = Lambda(s_k)
    sigma0 = (1/M) sum_k f_k               # normalizes the measure
    for n = 1 .. JMAX:                     # Taylor (Fourier/trapezoid) coeffs
        d_n = Re[ (1/M) sum_k f_k exp(-2*pi*i*(2n)*k/M) / R^(2n) ] / sigma0
    c[1] = d[1]                            # log-derivative recurrence
    for n = 2 .. JMAX:
        c[n] = d[n] - (1/n) sum_{i=1}^{n-1} i*c[i]*d[n-i]
    P[m] = (-1)^(m+1) * m * c[m]          # Hausdorff moments
    T(i,j) = P[i+j+1]                      # moment functional
    p[0] = 1 ;  norm[0] = <p0,p0>          # monic Gram-Schmidt
    for k = 1 .. N:
        x p[k-1]  ->  alpha[k-1] = <x p[k-1],p[k-1]>/<p[k-1],p[k-1]>
        orthogonalize against p[k-1] (alpha) and p[k-2] (b^2)
        b[k-2]^2  = norm[k-1]/norm[k-2]
        norm[k]   = <p[k],p[k]>
    assemble tridiagonal J_N = diag(alpha) + diag(b,+1) + diag(b,-1)
    verify  b_n^2 > 0  for all n           # positive-definiteness check
    lambda = descending eigenvalues of J_N ;  gamma = 1/sqrt(lambda)
    LOCK gamma if |gamma - gamma_zero|/gamma_zero < 1e-4
         against zeros found by an INDEPENDENT root finder (check only)
\end{verbatim}
\end{table*}

\noindent Settings: $\zeta$ and $\beta$ use $R=2$, $M=1536$, moments through
$P_{100}$ for $N=50$; the $\zeta\cdot\beta$ product uses the same $R,M$ with
$650$-digit precision and moments through $P_{210}$ for $N=100$; $\Delta$ uses
$R=4$, $M=1024$, $230$-digit precision, moments through $P_{211}$, $N=22$.
All linear algebra is standard double precision after the high-precision
coefficient extraction; eigenvalues are from dense symmetric tridiagonal
diagonalization. Independent zeros: $\zeta$ zeros from
\texttt{mpmath.zetazero}; $\beta$ zeros from Hurwitz-zeta root finding;
$\Delta$ zeros from $L(\Delta,s)$ root finding on the critical line.

\section{Experiment details and statistics}
\label{app:exp}

\paragraph{Experiment 1 (level repulsion).}
Locked levels: $\zeta$ $25$, $\beta$ $27$, $\Delta$ $15$ ($67$ levels,
$65$ adjacent spacings). Unfolding: cubic least-squares fit of the counting
function in $\gamma=1/\sqrt{\lambda}$ per family, normalized to mean gap
$1$; unfold residuals (max deviation of unfolded ranks from integers) are
$0.86$ ($\zeta$), $0.37$ ($\beta$), $0.35$ ($\Delta$). The minimum unfolded
gap is $0.432$ with zero gaps below $0.3$; a Poisson process of the same
length yields about $17$ such gaps ($20{,}000$ Monte-Carlo realizations,
$p<5\times10^{-5}$). The superposed $\zeta+\beta$ spectrum has minimum gap
$0.044$ and $7$ gaps below $0.3$.

\paragraph{Experiment 2 (dynamics).}
Time grids: $\tau\in[0.001,4.0]$ ($2000$ points) for $P(\tau)$ and
$\langle(\Delta n)^2\rangle$; $\tau\in[0.001,2.5]$ ($1000$ points) for
$K(\tau)$. The form factor is
$K(\tau)=N^{-1}\big|\sum_j e^{2\pi i\tau e_j}\big|^2$ on locked unfolded
levels; ramp slope is the least-squares slope on $[0.2,0.8]$. The GOE
benchmark is $200$ realizations of $60\times60$ symmetric Gaussian matrices
(seed $20260831$); the form factor uses the central $28$ levels (indices
$16$--$43$), dynamics the full spectrum. The uniform chain is
$\alpha_n=0,\ b_n=1$, $N=100$, with the exact unfolding $e_k=k$. Full-matrix
dynamics: $P(\tau)=|\sum_j w_j e^{-2\pi i\tau e_j}|^2$ with
$w_j=|\langle 0|j\rangle|^2$ from a boundary launch $|0\rangle$;
$\langle(\Delta n)^2\rangle=\sum_n n^2 p_n-(\sum_n n p_n)^2$. The locked
weight is the total $w_j$ on eigenstates whose $\gamma$ lies in the locked
range. For the product, the two families ($18$ $\zeta$, $35$ $\beta$ locked
levels) are unfolded separately and their form factors added incoherently,
$K=(N_\zeta K_\zeta+N_\beta K_\beta)/(N_\zeta+N_\beta)$.

\paragraph{Experiment 3 (evidence chain).}
\begin{itemize}[leftmargin=1.4em]
\item \emph{Coefficient tests} (on linearly detrended $\alpha_n$ and
$\log b_n$): autocorrelation $C(k)=\langle x_n x_{n+k}\rangle/\mathrm{Var}(x)$
to lag $20$, with a $95\%$ band from $300$ independent permutations;
power spectrum from the FFT of the standardized series, normalized to unit
mean; Wald--Wolfowitz runs test on the sign relative to the median
($z=(\text{runs}-\mu)/\sigma$; white noise gives $|z|\lesssim2$).
(Approximate entropy was computed but is not reported because, with only
$50$--$100$ samples, the zero-match frequency makes it short-sample
unstable; the runs and autocorrelation tests are the decisive
diagnostics.)
\item \emph{IPR}: $\mathrm{IPR}_k=\sum_n|\psi_{k,n}|^4/(\sum_n|\psi_{k,n}|^2)^2$.
The effective support is $N_{\rm eff}=\#\{n:b_n>0.05\,b_{\max}\}+2$
($22$ for $\zeta$, $11$ for $\beta$, $21$ for $\Delta$ at $N=50$/$22$);
locked-state IPR is computed on the effective support and the full-window
IPR over all sites. Anderson model: $H_{ii}$ uniform on $[-W/2,W/2]$,
unit hopping, $N=100$, $60$ realizations, central $50$ eigenstates.
\item \emph{Shuffle}: $80$ independent realizations; $\alpha_n$ and $b_n$
are permuted by two independent random permutations (marginal distributions
preserved exactly); $\langle r\rangle$ and IPR use the central half of the
full spectrum (all eigenvalues, physical finite system).
\item \emph{Perturbation}: $\alpha_n'=\alpha_n(1+\epsilon g_n)$,
$b_n'=b_n|1+\epsilon h_n|$ with Gaussian $g_n,h_n$; $40$ realizations per
$\epsilon\in\{10^{-3},3\cdot10^{-3},10^{-2},3\cdot10^{-2},10^{-1}\}$; a
level is tracked if it lies within $30\%$ relative $\gamma$ of the original
locked level.
\item \emph{Superposition}: locked $\gamma$ sets of two families are
concatenated and $\langle r\rangle$ computed directly (the gap ratio is
scale-free, so no unfolding is needed). Ensemble references (seed $7$,
$300$ realizations): GOE single $N=60$; two superposed GOE blocks
($40+40$ levels); two superposed Poisson sequences ($20+25$).
\item \emph{Finite size}: $\zeta,\beta$ at $N=20,25,\ldots,50$ (leading
principal submatrices of the $N=50$ coefficients); product at
$N=40,60,80,100$; $\Delta$ at $N=20,22$.
\end{itemize}
RNG seeds: $20260831$ (coefficient tests, IPR, shuffle, perturbation,
GOE dynamics) and $7$ (superposition Monte-Carlo).

\section{Jacobi coefficient tables (first ten sites)}
\label{app:coeffs}
On-site energies $\alpha_n$ and hopping amplitudes $b_n$; full tables
($N=50$ for $\zeta,\beta$, $N=100$ for the product, $N=22$ for $\Delta$)
are released as machine-readable files.

\begin{center}\footnotesize
\begin{tabular}{r|cc|cc|cc}
\toprule
$n$ & $\alpha_n^{(\zeta)}$ & $b_n^{(\zeta)}$
    & $\alpha_n^{(\beta)}$ & $b_n^{(\beta)}$
    & $\alpha_n^{(\Delta)}$ & $b_n^{(\Delta)}$\\
\midrule
0 & 0.001609 & 0.001911 & 0.011906 & 0.011934 & 0.00356 & 0.00448\\
1 & 0.003501 & 0.001149 & 0.017332 & 0.004843 & 0.00841 & 0.00258\\
2 & 0.001798 & 0.000697 & 0.007201 & 0.002721 & 0.00407 & 0.00164\\
3 & 0.001217 & 0.000554 & 0.004748 & 0.002117 & 0.00288 & 0.00121\\
4 & 0.001073 & 0.000498 & 0.003816 & 0.001591 & 0.00208 & 0.00090\\
5 & 0.000866 & 0.000353 & 0.002606 & 0.001074 & 0.00172 & 0.00084\\
6 & 0.000599 & 0.000272 & 0.001945 & 0.000955 & 0.00161 & 0.00073\\
7 & 0.000548 & 0.000287 & 0.001940 & 0.000913 & 0.00135 & 0.00065\\
8 & 0.000570 & 0.000256 & 0.001573 & 0.000657 & 0.00122 & 0.00053\\
9 & 0.000450 & 0.000208 & 0.001186 & 0.000575 & 0.00094 & 0.00043\\
\bottomrule
\end{tabular}
\end{center}

\section{Experimental prediction package (turnkey deliverable)}
\label{app:protocol}

This appendix is the self-contained interface for an experimental group.
Every number here is reproduced from the released data files and scripts.

\subsection{D.1 Hamiltonians to program}
Three one-dimensional nearest-neighbor tight-binding chains
$H_{nn'}=\alpha_n\delta_{nn'}+b_n(\delta_{n,n'+1}+\delta_{n+1,n'})$,
with coefficients taken verbatim from the released tables:
\begin{itemize}[leftmargin=1.4em]
\item $\zeta$ chain: $N=100$ (\texttt{zeta\_J100/hamiltonian.csv}), from the
higher-precision ($1300$-digit) construction: $b_n^2>0$ on all $99$
off-diagonals, $55$ locked zeros at relative error $\sim10^{-15}$; the earlier
$N=50$ table (\texttt{zeta\_J50\_matrix.csv}) is the same pipeline at lower
precision.
\item $\beta$ chain: $N=50$ (\texttt{beta\_J50\_matrix.csv}).
\item $\Delta$ chain: $N=22$ (\texttt{delta\_J22} coefficients, full
precision strings).
\end{itemize}
The hoppings $b_n$ decay super-exponentially; program the first $\sim30$
sites and truncate where $b_n$ meets the noise floor. Representative first
ten sites are in Appendix~C; the full tables are machine-readable in the
repository. The Heisenberg time is $\tau_H=1$ in the unfolded (mean-spacing)
convention used throughout.

\subsection{D.2 Observables and predicted values}
\begin{table*}[t]
\centering
\footnotesize
\caption{Observables and predicted values for the arithmetic chain versus controls.}
\label{tab:observables}
\begin{tabular}{lll}
\toprule
Observable & Arithmetic chain (prediction) & Controls \\
\midrule
Gap ratio $\langle r\rangle$ & $0.61$--$0.76$ (Wigner--Dyson) &
  P $0.386$; GOE $0.536$; GUE $0.600$\\
Near-degeneracies (gap$<0.3$, unfolded) & $\approx0$ in locked subspace &
  frequent for Poisson\\
Form factor $K(\tau)$ ramp slope & $0.74$ ($\zeta,\beta$), $0.74$ ($\Delta$) &
  GOE $0.79$; uniform chain $\approx0$\\
Revival at $\tau_H$ (boundary state) & absent; $P(\tau_H)=0.05$ (GOE-like) &
  uniform chain $0.99$\\
Eigenstate IPR (full window) & $0.056$ ($\zeta$), $0.060$ ($\beta$),
  $0.10$ ($\Delta$) & shuffled $0.49$--$0.58$\\
Coefficient runs test $|z|$ & $2.9$--$6.6$ (structured) & white noise $|z|\lesssim2$\\
\bottomrule
\end{tabular}
\end{table*}

\subsection{D.3 Control experiments}
\begin{enumerate}[leftmargin=1.4em]
\item \textbf{Uniform chain} $\alpha_n=0,\ b_n=1$: Poisson $\langle r\rangle$,
flat $K$, perfect revival---the integrable reference.
\item \textbf{Shuffled coefficients}: same $\{\alpha_n\},\{b_n\}$ values in a
random permutation (80 surrogates in numerics). Predicted: $\langle r\rangle\to
0.39$--$0.40$, IPR $\to0.49$--$0.58$. This is the decisive control isolating
arithmetic \emph{order} as the cause.
\item \textbf{Superposition}: merge the measured spectra of two independent
families ($\zeta+\beta$, $\beta+\Delta$). Predicted: $\langle r\rangle$
$0.50$--$0.52$, trending to Poisson; benchmark superposed-GOE $0.42$,
superposed-Poisson $0.38$.
\item \textbf{Genuine disorder (Anderson)}: replace $\alpha_n$ by i.i.d.
uniform $[-W/2,W/2]$. Predicted: extended for $W\lesssim2$, localized with
lost repulsion ($\langle r\rangle\to0.39$, IPR $0.39$ at $W=8$).
\item \textbf{Off-critical-line arithmetic control (Eisenstein $E_4$)}: the
Eisenstein series of weight~$4$ has $L(E_4,s)=\zeta(s)\zeta(s-3)$ with
elementary coefficients $\sigma_3(n)=\sum_{d|n}d^3$; its completed function
obeys $\Lambda(s)=\Lambda(4-s)$ about $s=2$, and its non-trivial zeros lie on
$\Re(s)=\tfrac12$ and $\Re(s)=\tfrac72$, i.e.\ \emph{offset} from the symmetry
center. Sampling the entire function $\Xi_{E_4}(s)=s(s-4)\Lambda(s)$ through
the identical pipeline ($650$-digit, $1536$ points) reproduces
$\Xi_{E_4}(2)=\tfrac1{72}$ to $\sim15$ digits, but the moment sequence changes
sign at order $8$ (predicted by the moment sum
$2\sum_k\Re(\gamma_k-\tfrac32 i)^{-2m}$, whose leading factor
$\cos(2m\theta_1)$ with $\theta_1=\arctan(3/2\gamma_1)\approx0.106$ first turns
negative at $m\approx7.4$), and the Jacobi chain fails at order~$4$ ($b_n^2$
non-positive; signs $+--+--++-\cdots$). The long $b_n^2>0$ chains of
$\zeta,\beta,\Delta$ therefore track the zeros sitting on the symmetry line,
rather than being a generic output of the numerics: the pipeline discriminates
on-line from off-line zero distributions. (\texttt{controls/e4\_eisenstein/}.)
\end{enumerate}

\subsection{D.4 Hardware noise budget}
With multiplicative calibration noise of relative size $\epsilon$
($a_n\to a_n(1+\epsilon g_n)$, $b_n\to b_n|1+\epsilon h_n|$, $g,h$ standard
normal):
\begin{itemize}[leftmargin=1.4em]
\item $\epsilon\le0.1$: statistics unbiased, every level trackable---\textbf{safe
operating regime}.
\item $\epsilon_c\approx0.2$--$0.3$: crossover midpoint ($\langle r\rangle
\approx0.50$).
\item $\epsilon\ge0.5$: Poisson floor $0.41$; positive-eigenvalue fraction
falls below $0.82$ (structural damage).
\end{itemize}
\textbf{Engineering requirement:} calibrate on-site frequencies and couplings
to $\lesssim10\%$ relative error; $1$--$3\%$ gives a pristine ramp.

\subsection{D.5 Pass/fail judgment}
The arithmetic-origin mechanism is \textbf{observed} only if, in the
\emph{same} device in the \emph{same} calibration run, the ordered $\zeta$
(or $\beta$) chain shows (i) $\langle r\rangle$ in the Wigner band with no
resolved near-degeneracies, (ii) a linear $K(\tau)$ ramp with plateau,
(iii) no $\tau_H$ revival, and (iv) low-IPR extended states,
\textbf{and the shuffled-coefficient control simultaneously shows Poisson
statistics together with localization}, while the uniform chain shows
Poisson statistics and its perfect revival. The controls are not optional
corroboration: they are the null test. \textbf{If the shuffled control
still shows Wigner--Dyson statistics on the same device, the
arithmetic-origin claim is falsified---even if the ordered chain looks
chaotic.} Likewise, a chaotic-looking ordered chain without the same-device
controls does not count as evidence.

\begin{acknowledgments}
We thank the open-source mathematical computing community for the tools that
made this work possible: mpmath for high-precision arithmetic, NumPy/SciPy for
linear algebra, and Matplotlib for visualization. All computations were
performed on cloud computing resources. The author acknowledges helpful
discussions with AI-assisted research collaborators that contributed to the
numerical experiments and code development.
\end{acknowledgments}

\begin{widetext}
\textbf{Data availability.} All code, data, and reproducible scripts are
publicly available on GitHub at
\href{https://github.com/maomao8412/hilbert-polya-jacobi}{github.com/maomao8412/hilbert-polya-jacobi}.
Archived versions are deposited on Zenodo (DOIs: 10.5281/zenodo.22180921,
10.5281/zenodo.22181443, 10.5281/zenodo.22194207).
\end{widetext}

\begin{thebibliography}{9}
\bibitem{berry1977} M.~V. Berry and M.~Tabor,
\emph{Level clustering in the regular spectrum}, Proc. R. Soc. Lond. A
\textbf{356}, 375 (1977).
\bibitem{bgs1984} O.~Bohigas, M.-J. Giannoni, C.~Schmit,
\emph{Characterization of chaotic quantum spectra and universality of level
fluctuation laws}, Phys. Rev. Lett. \textbf{52}, 1 (1984).
\bibitem{mehta} M.~L. Mehta, \emph{Random Matrices}, Academic Press (2004).
\bibitem{atas2013} Y.~Y. Atas, E.~Bogomolny, O.~Giraud, G.~Roux,
\emph{Distribution of the ratio of consecutive level spacings in random
matrix ensembles}, Phys. Rev. Lett. \textbf{110}, 084101 (2013).
\bibitem{montgomery1973} H.~L. Montgomery, \emph{The pair correlation of
zeros of the zeta function}, Proc. Symp. Pure Math. \textbf{24}, 181
(1973).
\bibitem{odlyzko1987} A.~M. Odlyzko, \emph{On the distribution of spacings
between zeros of the zeta function}, Math. Comp. \textbf{48}, 273 (1987).
\bibitem{rudnick1996} Z.~Rudnick and P.~Sarnak, \emph{Zeros of principal
$L$-functions and random matrix theory}, Duke Math. J. \textbf{81}, 269
(1996).
\bibitem{sarnak1995} P.~Sarnak, \emph{Arithmetic quantum chaos}, Israel
Math. Conf. Proc. \textbf{8}, 183 (1995).
\bibitem{anderson1958} P.~W. Anderson, \emph{Absence of diffusion in certain
random lattices}, Phys. Rev. \textbf{109}, 1492 (1958).
\bibitem{neill2016} C.~Neill, P.~Roushan, M.~Fang et al.,
\emph{Ergodic dynamics and thermalization in an isolated quantum system},
Nat. Phys. \textbf{12}, 1037 (2016).
\bibitem{kaufman2016} A.~M. Kaufman, M.~E. Tai, A.~Lukin, M.~Rispoli,
R.~Schittko, P.~M. Preiss, M.~Greiner, \emph{Quantum thermalization through
entanglement in an isolated many-body system}, Science \textbf{353}, 794
(2016).

\bibitem{maldacena2016} J.~Maldacena, S.~H. Shenker, D.~Stanford,
  \emph{A bound on chaos}, JHEP \textbf{08} (2016) 106.

\bibitem{chen2026rh} Z.~Chen, \emph{The Riemann hypothesis via Hankel
positivity of the $\xi$-function log-moments},
Zenodo doi:10.5281/zenodo.22167882 (2026).
\bibitem{chen2026grh} Z.~Chen, \emph{Generalized Riemann hypothesis for
Dirichlet $L$-functions via three-term polylogarithm decomposition},
Zenodo doi:10.5281/zenodo.22167742 (2026).
\bibitem{berry1999} M.~V. Berry and J.~P. Keating,
  \emph{The Riemann zeros and eigenvalue asymptotics}, SIAM Review \textbf{41}, 236 (1999).

\bibitem{connes1999} A.~Connes,
  \emph{Trace formula in noncommutative geometry and the zeros of the Riemann zeta function},
  Selecta Math. (N.S.) \textbf{5}, 29 (1999).

\bibitem{suo2025} X.~Suo et al.,
  \emph{On the Hamiltonian with energy levels corresponding to Riemann zeros},
  arXiv:2505.21192 (2025).

\end{thebibliography}

\end{document}
